\chapter{非交换代数.}

我们已经看到（第 62 页），第一批非交换代数如何于 1843-44 年出现在哈密顿 {[}145 a{]}和格拉斯曼（{[}134{]}，第 I₂ 卷）的工作中。哈密顿在引入四元数时，就已经对实数域上任意有限秩的代数有了非常清楚的概念（{[}145 a{]}，序言，第 (26)-(31) 页）。¹在发展他的理论时，他稍后有了考虑他所谓“双四元数”的想法，也就是说复数域上与四元数体有相同乘法表的代数；他借这个机会观察到，这个扩张所产生的效果是零因子的出现（{[}145 a{]}，第 650 页）。格拉斯曼的观点有些不同，而且他的“外代数”将在很长一段时间里或多或少独立于代数理论之外，²但用他那仍然缺乏精确性的语言，人们不可能认不出由生成元和关系组所定义的（有限维或非有限维）代数的最初思想（{[}134{]}，第 II₁ 卷，第 199-217 页）。

1850-1860 年间，一些新的代数例子以或多或少明确的方式被引入：如果说凯莱在发展矩阵理论时（{[}58{]}，第 II 卷，第 475-496 页）还没有考虑方阵构成一个代数这件事（这个观点直到 1870 年前后才由皮尔斯父子明确表达 {[}248 c{]}），那么他至少已经在这个场合注意到，存在一个 2 阶矩阵组满足四元数的乘法表，这条评注可以看作一个代数的线性表示的第一个例子。\(^{3}\)另一方面，在他定义有限群的抽象概念的那部著作中，他还顺带给出了这样一个群的代数的定义，却没有从这个定义中得出任何东西（{[}58{]}，第 II 卷，第 129 页）。

在 1870 年之前没有其他值得注意的进展可以指出；但在那个时候，关于（实数域或复数域上）有限维代数的一般结构的研究开始了。B. 皮尔斯沿这条路迈出了最初的几步；他引入了幂零元素、幂等元素的概念，证明了一个代数（有或没有单位元）只要至少有一个元素不是幂零的，就包含一个幂等元 \(\neq 0\)，并写下了著名的分解

\[
\boldsymbol {x} = e x e + (\boldsymbol {x} e - e x e) + (e \boldsymbol {x} - e x e) + (\boldsymbol {x} - \boldsymbol {x} e - e \boldsymbol {x} + e x e)
\]

（e 为幂等元，x 为任意元素），并且有了把一个幂等元分解为两两正交的“本原”幂等元之和的想法（尚有些不精确）{[}247{]}。同样，据克利福德所说（{[}65{]}，第 274 页），\(^{4}\)两个代数的张量积的概念也应归功于 B. 皮尔斯，克利福德本人把它隐含地应用于哈密顿“双四元数”的一个推广（{[}65{]}，第 181-200 页），并在几年后明确地应用于冠有其名的那些代数的研究（{[}65{]}，第 397-401 页和第 266-276 页）。B. 皮尔斯把这些新概念用于（复数域上）小维数代数的分类，这个问题在 1880 年前后也被英美学派的其他数学家、以凯莱和西尔维斯特为首的数学家们所攻克。可能的结构的一大堆花样由此迅速变得明显，而无疑正是这个事实，在随后的时期里把研究引向获取具有更特殊性质的代数类。

在欧洲大陆，观念的演变相当不同，这类研究在 1880 年之前就已出现。1878 年，弗罗贝尼乌斯证明四元数构成实数域上（有限维）除环的唯一例子（{[}119{]}，第 I 卷，第 343-405 页）——这个结果两年后由 C. S. 皮尔斯独立发表 {[}248 d{]}。从 1861 年起，魏尔斯特拉斯在准确陈述高斯的一条评注时，就在他的讲课中把 R 或 C 上没有幂零元素的交换代数\(^{5}\)刻画为域（与 R 或 C 同构的域）的直和；戴德金则在他那一方面，大约在 1870 年，在联系于他对交换域理论的“超复”概念时得到了相同的结论；他们的证明发表于 1884-85 年（{[}329 a{]}，第 II 卷，第 311-332 页和 {[}79{]}，第 II 卷，第 1-19 页）。也是在 1884 年，H. 庞加莱在一篇很短、非常隐晦的注记中（{[}251 a{]}，第 V 卷，第 77-79 页），提请人们注意可以把方程 \(z_{i} = \phi_{i}(x_{1}, \ldots, x_{n}, y_{1}, \ldots, y_{n})\)——它们表示一个代数中的乘法法则 \((\sum_{i} x_{i} e_{i})(\sum_{i} y_{i} e_{i}) = \sum_{i} z_{i} e_{i}\)——看作（当然是局部地）定义一个李群。这条评注似乎给李及其学生们（施图迪、舍弗斯、F. 舒尔，稍后还有莫利安和 E. 嘉当）留下了深刻印象，他们此时正忙于发展“连续”群的理论，特别是分类问题（特别参见 {[}271{]}，第 387 页）；在 1885-1905 年间，它引导这个学派的数学家把与他们在李群和李代数研究中使用的同类方法应用于代数结构的研究。这些方法首先依赖于对一个元素相对于其正则表示的特征多项式（这个多项式在前面引用的魏尔斯特拉斯和戴德金的工作中已经遇到过）的考虑，以及对这个多项式分解为不可约因子的考虑；在这种分解中，正如弗罗贝尼乌斯稍后将发现的，反映着正则表示分解为不可约成分。

在李学派关于代数的研究过程中，理论的“内蕴”概念逐渐显露出来。根基的概念以一种特殊的情形（商掉根数为域的直和的情形）出现于 G. 舍弗斯 1891 年的工作 {[}271{]}，在莫利安 {[}224 a{]}和嘉当（{[}52 a{]}，第 II \(_{1}\) 卷，第 7-105 页）那里则更为清晰，他们研究的是一般情形（“根基”一词本身出自弗罗贝尼乌斯（{[}119{]}第 III 卷，第 284-329 页））。施图迪和舍弗斯 {[}271{]}突出了一个代数作为若干个别的代数之直和的概念（已由 B. 皮尔斯预见到（{[}247{]}，第 221 页））。最后，在莫利安 {[}224 a{]}那里引入了一个代数的商代数，这个概念本质上等价于双边理想（第一次由嘉当定义（{[}52 a{]}，第 II \(_{1}\) 卷，第 7-105 页））或同态（这个名称也应归功于弗罗贝尼乌斯）的概念；这里与群的类比非常鲜明，稍后的 1904 年，爱普斯坦和韦德伯恩考虑了双边理想的合成列，并把若尔当-赫尔德定理推广到它们。这一时期最重要的结果是 T. 莫利安 {[}224 a{]}的结果：在单群概念的指引下，他定义了（C 上的）单代数并证明这些代数是矩阵代数，然后证明任意一个有限秩代数在基域

C 上的结构本质上归结为商掉根数为域的直和的情形（舍弗斯已研究过的情形）。这些结果不久之后就被 E. 嘉当以一种更清楚、更严格的方式重新找到并建立（{[}52 a{]}，第 II \(_{1}\) 卷，第 7-105 页），他借这个机会引入了半单代数的概念，并提出了与域 C 上任意代数相联系的一些数值不变量（“嘉当整数”）——从而把这个理论带到了一个我们此后几乎没有超越的地步；\(^{6}\)最后他把莫利安的结果和他自己的结果推广到 R 上的代数。

大约在 1900 年前后，兴起了一场观念运动，它导致在与线性代数有关的一切事务中放弃对基域的所有限制；特别地，我们必须指出以 E. H. 穆尔和 L. E. 迪克森为核心的美国学派对有限域研究所给予的强大推动；这项研究最引人注目的结果是韦德伯恩定理 {[}328 a{]}，它证明每个有限域都是交换的。1907 年，韦德伯恩重新接过嘉当的结果并把它们推广到任意基域 {[}328 b{]}；在这样做时，他完全抛弃了他的先辈们的方法（一旦基域不再是代数闭的或全序的，这些方法就不再适用），并回到——同时加以完善——B. 皮尔斯的幂等元技巧，这使他能够把关于半单代数结构（其研究被归结为除环的研究）的定理置于定论的形式。此外，基域的扩张问题在他所采取的视角下被自然地陈述，他证明了每个半单代数在基域的一个可分扩张之后仍是半单的，\(^{7}\)而且如果这个扩张取得足够大，它就变成中心矩阵代数的直和（{[}328 b{]}，第 102 页）。\(^{8}\)稍后，迪克森就 \(n = 3\)（{[}88 a{]}）而韦德伯恩本人就任意 \(n\)（{[}328 c{]}）给出了在其中心上秩为 \(n^{2}\) 的除环的头几个例子，\(^{9}\)从而在一个特殊情形中开创了后来由 R. 布饶尔 {[}34 a{]}和 E. 诺特 {[}236 c{]}加以发展的“交叉积”和“因子组”的理论。最后，1921 年，韦德伯恩证明了交换定理的一个特殊情形 {[}328 d{]}。

与此同时，从 1896 年到 1910 年，在弗罗贝尼乌斯、伯恩赛德和 I. 舒尔手中，发展起一个与代数理论邻近的理论，即群的线性表示理论（起初限于有限群的表示）。它源自戴德金的一些评注：后者（甚至在他关于代数的著作发表之前）大约在 1880 年，在研究伽罗瓦扩张的正规基时遇到了“群行列式”\(\det(x_{st-1})\)，其中 \((x_s)_{s \in G}\)是以有限群 \(G\) 为指标集的一组未定元（换句话说，即群代数 \(G\) 的泛元素相对于其正则表示的范数）；他观察到，当 \(G\) 交换时，这个多项式分解为线性因子（这推广了很久以前就“循环”行列式证明的一个恒等式，那些行列式对应于循环群 \(G\)）。在他与弗罗贝尼乌斯的非常有趣的通信（{[}79{]}，第 II 卷，第 414-442 页）中，戴德金于 1896 年提请他注意这一性质、它与交换群特征标理论 {[}327 c{]}的联系，以及他在 1886 年得到的关于一些特殊非交换群的类似结果。几个月后，弗罗贝尼乌斯完全解决了“群行列式”分解为不可约因子的问题（{[}119{]}，第 III 卷，第 38-77 页），这要归功于他对特征标概念的光辉推广（{[}119{]}，第 III 卷，第 1-37 页），我们不应当在此谈论它。但必须指出，在这个理论后来的发展中，\(^{10}\)弗罗贝尼乌斯始终清楚它与代数理论的联系（戴德金在他的信中也没有停止强调这一点）；在为群引入了不可约表示和完全可约表示的概念（{[}119{]}，第 III 卷，第 82-103 页）、并证明了正则表示包含全部不可约表示之后，正是用类似的方法，他于 1903 年提出重新接过莫利安-嘉当的理论（{[}119{]}，第 III 卷，第 284-329 页）。在伯恩赛德 {[}44 a{]}和 I. 舒尔 {[}279 c{]}那里，理论的“超复”方面没有明确出现；但正是与他们一起，不可约表示的基本性质、舒尔引理和伯恩赛德定理得以显露。最后，就我们的目的而言必须指出，正是在这个理论中第一次出现了交换定理的两个特殊情形：一是在 I. 舒尔的学位论文 {[}279 a{]}中，他把线性群的表示与对称群的表示联系起来（恰恰是通过张量空间自同态环中的交换化）；二是在他 1905 年的工作 {[}279 c{]}中，他证明与域 \(\mathbf{C}\) 上一个不可约表示的所有矩阵都可交换的矩阵都是 \(I\) 的标量倍（这个结果也可以由伯恩赛德定理得出）。

还剩下的是把这些理论的共同基底清楚地剥离出来：这是 E. 诺特和 E. 阿廷周围的德国学派在 1921-1933 年间的工作，这个时期见证了现代代数的创立。早在 1903 年，H. 庞加莱在一部关于线性微分方程的代数积分的著作中（{[}251 a{]}，第 III 卷，第 140-149 页），就在一个代数中定义了左理想和右理想以及极小理想的概念；他还注意到，在一个半单代数中，每个左理想都是它与各单分支的交的直和，而且在 n 阶矩阵代数中，极小理想的维数是 n；但他的工作没有引起代数学家的注意。\(^{11}\)1907 年，韦德伯恩重新定义了一个代数的左理想和右理想，并证明了若干性质（特别是根是最大的幂零左理想这一点（{[}328 b{]}，第 113-114 页））。但我们必须等到 1927 年，这些概念才在代数理论中得到本质性的使用。\(^{12}\)把此前在各处零星出现的证明程序置于一般形式，\(^{13}\)W. 克鲁尔在 1925 年 {[}187 a{]}和 E. 诺特在 1926 年 {[}236 b{]}引入并系统地使用了极大条件和极小条件；前者用它把雷马克关于一个有限群分解为不可分解群之直积的定理推广到带算子的交换群（他在这个场合定义了这种群），而后者把这些条件用于戴德金环的刻画。1927 年，E. 阿廷 {[}7 c{]}把同一思想应用于非交换环，表明如何通过对极小理想的系统研究，把韦德伯恩的定理推广到其左理想同时满足极大条件和极小条件的所有环。\(^{14}\)

另一方面，克鲁尔在 1926 年 {[}187 b{]}把带算子交换群的概念与群的线性表示的概念联系了起来；这个观点被 E. 诺特在 1929 年的一部基本著作 {[}236 c{]}中推广到代数并得到详细发展，这部著作无论就所引入思想的重要性还是就阐述的清晰而言，都值得与施泰尼茨关于交换域的著作并列为现代线性代数的支柱之一。\(^{15}\)

最后，在始于 1927 年的一系列工作中（{[}237{]}、{[}34 a{]}、{[}236 d{]}），E. 诺特和 R. 布饶尔（从 1929-31 年起 A. 阿尔伯特和 H. 哈塞也加入进来）在韦德伯恩和迪克森停下的地方重新接过左除环的研究。如果说他们的结果中最重要的部分在于对布饶尔群的深入研究（特别是在代数数域上），因而超出了本注记的范围，那么无论如何要指出，正是在这项工作的过程中，交换定理被精确地陈述，还有中性化一个单代数的域的概念及其与极大交换子域的关系；最后，1927 年，斯科朗刻画了单环的自同构 {[}286 b{]}，这个定理几年后被 E. 诺特 {[}236 b{]}和 R. 布饶尔 {[}34 a{]}重新发现。

于是，到 1934 年，单环和半单环的初等理论差不多已达到其最终面貌（关于理论在这一阶段的状态的综述，见 {[}87{]}）；从那时起，它朝两个不同的方向发展，我们只限于简要地提及。一方面，R. 布饶尔和 E. 诺特的“因子组”理论近来在它被并入现代同调代数之后获得了新的推动。\(^{16}\)另一方面，人们作过很多尝试——成败不一——把经典理论的结果至少部分地推广到没有极小条件的环\(^{17}\)或没有单位元的环。但迄今为止，这些推广对数学的其他分支几乎没有产生什么影响；关于这些工作的更多细节，我们请读者参看 N. 雅各布森最近的综述 {[}172 b{]}。
% semantic-fragments: legacy-input-seam
\relax{}
